\section{Assembly of a smooth metric germ}\label{sec:assembly}

We now turn the whole-patch obstruction into one metric germ.
There are two levels of accumulation: obstructed patches approach
every boundary point of each negative disk, and the disks approach
the origin. Summable bounds in each derivative order will preserve
smoothness at both levels.
We first recall the elementary prescribed-curvature construction
used in \cite[p.~453]{Khuri2009}.

\begin{lemma}[Prescribing small curvature]\label{lem:prescribe}
Let \(K_0\in C^\infty((-1,1)^2)\) with
\(\norm{K_0}_\infty\le10^{-3}\).
There is a smooth positive metric
\[
\bar g=\dd X^2+f(X,Y)^2\dd Y^2
\]
whose curvature is \(K_0\). Its eigenvalues have positive uniform
lower and upper bounds.
\end{lemma}

\begin{proof}
Solve the ordinary differential equation with parameter \(Y\),
\begin{equation}\label{eq:prescribe-ode}
f_{XX}=-K_0f,\qquad f(0,Y)=1,\qquad f_X(0,Y)=0.
\end{equation}
Smooth parameter dependence gives smoothness. On any segment
joining \(0\) to \(X\), let \(F\) be the supremum of \(\abs f\).
The integral equation gives
\[
\sup\abs{f-1}\le10^{-3}F,\qquad F\le1+10^{-3}F.
\]
Thus \(f\) stays uniformly close to one and, in particular,
strictly positive. Computing the Christoffel symbols of
\(\dd X^2+f^2\dd Y^2\) gives \(K_{\bar g}=-f_{XX}/f\).
This formula holds also when \(f\) depends on \(Y\).
Equation \eqref{eq:prescribe-ode} proves the assertion.
\end{proof}

\subsection{Placement and smooth summation}

Use the fixed coordinates \(X,Y\) on \((-1,1)^2\).
Let \(\mathcal D_j\), \(j\ge1\), be the round disks with centers
\((2^{-j},0)\) and radii \(2^{-j-4}\).
They have disjoint closures and disjoint slightly larger
neighborhoods whose closures avoid the origin, and converge to the
origin without containing it.

Choose a finite set of rotations whose first-axis lines form an
angular net of mesh less than \(1/1000\).
Place countably many affine patch maps
\[
\Phi_n(x,y)=c_n+\lambda_nR_n(x,y),\qquad \lambda_n>0,
\]
where \(R_n\) belongs to this finite set, with the following properties:
\begin{enumerate}[label=(\roman*),leftmargin=2.3em]
\item the images of \([-4,4]^2\) have disjoint closures and lie
outside all the closed disks;
\item for every \(j\), every \(b\in\partial\mathcal D_j\), and every
rotation in the net, a sequence of full patches with that rotation
converges to \(b\), with diameters tending to zero.
\end{enumerate}
To obtain this placement, use disjoint thin annuli outside each
\(\mathcal D_j\), within its larger neighborhood, converging to
its boundary. In each annulus choose finitely many distinct centers
so that those assigned each rotation form a successively finer
angular net on the circle. Their number is finite, so all square
sizes can be chosen small enough to give disjoint closed patches
inside that annulus. This proves both properties. \Cref{fig:accumulation} separates the
two levels of the construction.

\begin{figure}[H]
\centering
\begin{tikzpicture}[font=\small,>=Latex]
  \begin{scope}
    \draw[->,gray] (-.2,0)--(4.1,0) node[right,black] {\(X\)};
    \fill (0,0) circle (1.5pt) node[below=4pt] {\(0\)};
    \foreach \position/\radius/\disklabel in {3.2/.20/1,1.6/.10/2,.8/.05/3} {
      \draw[fill=orange!18] (\position,0) circle (\radius);
      \node[above=10pt] at (\position,0) {\(\mathcal D_{\disklabel}\)};
    }
    \fill (.4,0) circle (.025);
    \fill (.2,0) circle (.0125);
    \draw[->,blue!65!black] (2.6,-.50)--(.6,-.50);
    \node at (2,-1.05) {(a) Disks approaching the origin};
  \end{scope}
  \begin{scope}[shift={(7.1,.18)}]
    \draw[fill=orange!18] (0,0) circle (.67);
    \draw[gray,dashed] (0,0) circle (.90);
    \draw[gray,dashed] (0,0) circle (1.23);
    \foreach \angle/\innerrotation/\outerrotation in {
      10/0/60,50/30/0,90/60/30,130/15/75,170/45/15,
      210/75/45,250/0/30,290/60/0,330/30/60} {
      \begin{scope}[shift={(\angle:.923)},rotate=\innerrotation]
        \draw[blue!65!black,fill=blue!8] (-.033,-.033) rectangle (.033,.033);
      \end{scope}
      \begin{scope}[shift={(\angle:1.25)},rotate=\outerrotation]
        \draw[blue!65!black,fill=blue!8] (-.07,-.07) rectangle (.07,.07);
      \end{scope}
    }
    \node at (0,0) {\(\mathcal D_j\)};
    \node at (0,-1.55) {(b) Patches approaching a disk};
  \end{scope}
\end{tikzpicture}
\caption{The two accumulation levels, schematically and not to scale.
The disks in (a) are disjoint and avoid the origin. In (b), disjoint
full patches occupy successive annuli outside one disk. Only a few
patches and annuli are shown. Patch orientations are independent of
their positions around the disk; the actual construction approaches
every boundary point in every direction of a fixed finite rotation net.}
\label{fig:accumulation}
\end{figure}


Now prescribe \(K_0\). In each \(\mathcal D_j\) choose a smooth
negative bump, strictly negative throughout the disk and flat on
its boundary. In each patch choose a smooth bump supported in
\(\Phi_n((-4,4)^2)\), equal on a neighborhood of
\(\Phi_n([-3,3]^2)\) to
\begin{equation}\label{eq:patch-K0}
K_0(\Phi_n(x,y))=\kappa_n(x^2-h(y)),\qquad \kappa_n>0.
\end{equation}
The amplitudes of all disk and patch bumps can be chosen so that
their sum is smooth and \(\norm{K_0}_\infty\le10^{-3}\).
Here is a precise choice. Enumerate the unscaled, globally smooth
bump profiles, extended by zero, as \(k_n\). Choose positive
numbers \(\epsilon_n\) such that
\[
\norm{\epsilon_nk_n}_{C^m}\le10^{-3}2^{-n}
\qquad (0\le m\le n).
\]
There are only finitely many restrictions at stage \(n\).
For each fixed \(m\), the resulting series of derivatives has
a uniformly summable tail. Its sum is \(C^\infty\), including
on all accumulation sets. Disjointness of the supports preserves
the strict negativity inside each disk and the exact profiles
\eqref{eq:patch-K0}. Every summand is flat at each disk-boundary
point at which its support can accumulate, so
\[
K_0=0\qquad\text{on }\partial\mathcal D_j
\]
for every \(j\).

Apply \Cref{lem:prescribe} to obtain \(\bar g\).
For each patch, the actual pullback \(g_{0,n}=\Phi_n^*\bar g\)
has curvature
\[
K_{g_{0,n}}=K_{\bar g}\circ\Phi_n
=\kappa_n(x^2-h(y)).
\]
Indeed, in coefficient matrices,
\(g_{0,n}=\lambda_n^2R_n^T(\bar g\circ\Phi_n)R_n\);
curvature is a scalar under pullback.
This is a background metric of the kind required in
\Cref{sec:model}.

By \Cref{prop:patch}, choose a smooth tensor perturbation
\(\eta_n\) supported in the model box \(\mathcal P\) such that
\(g_{0,n}+\eta_n\) admits no admissible height.
Push \(\eta_n\) to the original coordinates and extend by zero,
obtaining \(\widehat\eta_n\).
The extension is smooth because its support is strictly inside
the full patch. In matrix notation the pushforward is
\[
\widehat\eta_n
=\lambda_n^{-2}R_n(\eta_n\circ\Phi_n^{-1})R_n^T.
\]
Consequently, for each ordinary derivative order \(m\),
\[
\norm{D^m\widehat\eta_n}_\infty
\le C_m\lambda_n^{-m-2}\norm{D^m\eta_n}_\infty,
\]
where \(C_m\) accounts only for the fixed-dimensional component
norms and rotation. The perturbation may be chosen arbitrarily
small in every prescribed finite collection of seminorms, so
these fixed scaling costs can be paid separately on each patch.
Require
\begin{equation}\label{eq:tensor-sum}
\norm{\widehat\eta_n}_{C^m}\le\epsilon\,2^{-n}
\qquad(0\le m\le n),
\end{equation}
where \(\epsilon>0\) is smaller than half a uniform lower
eigenvalue bound for \(\bar g\).
The metric
\begin{equation}\label{eq:final-metric}
g=\bar g+\sum_{n=1}^\infty\widehat\eta_n
\end{equation}
is therefore smooth and positive.
It retains the chosen patch obstruction on every patch:
all the other tensor perturbations vanish there.
Every tensor perturbation and all its derivatives vanish on
each closed disk. Uniform convergence of the derivative series
in \eqref{eq:tensor-sum} consequently preserves the background
metric jets on those disks and their boundaries. Thus
\begin{equation}\label{eq:disk-curvature}
K_g<0\text{ in }\mathcal D_j,\qquad
K_g=0\text{ on }\partial\mathcal D_j.
\end{equation}
This construction imposes no lower bound on the individual
\(\kappa_n\), \(\lambda_n\), or admissible Hessian separations.

\subsection{The boundary obstruction}

The following graph observation is also used in the geometric
construction of \cite[Section~2, Step~3]{NY2002}. We include its short proof.

\begin{lemma}[A negatively curved disk]\label{lem:disk}
Suppose a smooth metric on a neighborhood of a closed disk
has negative curvature in the disk and zero curvature on its
boundary. If an isometric immersion of this neighborhood is
a graph over a plane, its second fundamental form is nonzero
at some boundary point. At that point it has rank one.
\end{lemma}

\begin{proof}
Let \(u\) be the graph function and let \(\Omega\) be the planar
projection of the disk. This is a bounded domain with connected
smooth boundary, and \(u\) is smooth on a neighborhood of
\(\overline\Omega\). Suppose the second fundamental form vanishes
at every boundary point. The graph formula
\[
\II_{ij}=\frac{u_{ij}}{\sqrt{1+\abs{Du}^2}}
\]
then gives \(D^2u=0\) along \(\partial\Omega\). Differentiation
along this connected curve shows that \(Du\) is a fixed vector
\(a\) there. Hence \(u=a\cdot x+c\) on the boundary for a constant
\(c\).

The function \(v=u-a\cdot x-c\) is zero on \(\partial\Omega\)
and has
\[
\det D^2v=\det D^2u<0\qquad\text{in }\Omega.
\]
It cannot vanish identically. If it is nonzero, compactness
gives either a positive interior maximum or a negative interior
minimum. Its Hessian is semidefinite there, with nonnegative
determinant, a contradiction.
Thus the second fundamental form is nonzero at some boundary
point. The Gauss equation and zero boundary curvature make
its determinant zero, so its rank is one.
\end{proof}

\begin{proof}[Proof of Theorem~A]
Use the metric \eqref{eq:final-metric}.
Suppose that a neighborhood of the origin admits a smooth
isometric immersion \(F\) into \(\R^3\).
Projection onto the tangent plane at \(F(0)\) is a local
diffeomorphism. After restricting the domain, the immersion
is a graph and the projection is a diffeomorphism onto a
planar open set. Choose \(j\) so large that
\(\overline{\mathcal D_j}\) lies in this restricted domain.
By \eqref{eq:disk-curvature} and \Cref{lem:disk}, there is
\(b\in\partial\mathcal D_j\) where \(\II\) is nonzero and
has rank one.

Let \(e\) be the fixed unit normal at \(F(b)\), and set
\(z=\langle F,e\rangle\). By \Cref{lem:height}, the covariant
Hessian \(H\) is nonzero and rank one at \(b\), while \(dz(b)=0\).
In particular \(1-\abs{dz}_g^2>0\) nearby.

Regard \(H(b)\) as a symmetric form in the fixed \(X,Y\)
parameter plane. Choose a rotation in our finite net whose
first-axis line makes angle less than \(1/1000\) with its
kernel. In that rotated parameter frame, a rank-one form
has
\[
H_{yy}(b)\ne0,\qquad
\abs{H_{xy}(b)/H_{yy}(b)}
=\abs{\tan\alpha}<1/100,
\]
where \(\alpha\) is the angular error.
Continuity gives the same strict inequalities in a
neighborhood of \(b\).

There are full patches of this orientation converging to \(b\).
One such patch lies wholly in that neighborhood and in the
domain of \(F\). Its common dilation of the two axes multiplies
all Hessian entries by the same positive factor, leaving their
ratio unchanged. In model coordinates the pulled-back height
is therefore smooth on a neighborhood of \(S\), has \(E>0\),
\(H_{yy}\ne0\), and \(\abs q\le1/100\) throughout \(S\).
It satisfies the Darboux equation by \Cref{lem:height}.
It is an admissible height for a metric chosen to have none,
contradicting \Cref{prop:patch}.

Thus no neighborhood of the origin admits the asserted smooth
isometric immersion. The open square with the smooth positive
metric \(g\) is the required Riemannian surface.
\end{proof}

\begin{corollary}[Euclidean Taylor jet]\label{cor:flat-jet}
The metric in Theorem~A can be chosen so that, in the coordinates
\(X,Y\),
\[
\partial^\alpha(g_{ij}-\delta_{ij})(0)=0
\qquad
\text{for every multi-index }\alpha\in\mathbb N_0^2
\text{ and }i,j\in\{1,2\}.
\]
Thus smooth local isometric realizability is not determined by
the full Taylor jet of the metric at the base point.
\end{corollary}

\begin{proof}
Every individual curvature bump and every tensor
\(\widehat\eta_n\) vanishes on a neighborhood of the origin.
The uniform convergence of each derivative series therefore gives
\[
\partial^\alpha K_0(0)=0,
\qquad
\partial^\alpha\Bigl(\sum_n\widehat\eta_n\Bigr)(0)=0
\]
for every multi-index \(\alpha\).
It remains to check the background metric
\(\bar g=\dd X^2+f^2\dd Y^2\).
The initial conditions in \eqref{eq:prescribe-ode} give
\(\partial_X^a\partial_Y^b(f-1)(0)=0\) for \(a=0,1\)
and every \(b\ge0\). For \(a\ge2\), the same equation yields
\[
\partial_X^a\partial_Y^b(f-1)(0)
=-\partial_X^{a-2}\partial_Y^b(K_0f)(0)=0,
\]
since every term in the Leibniz expansion contains a derivative
of \(K_0\) at the origin. Hence \(f-1\), and therefore
\(\bar g_{ij}-\delta_{ij}\), has zero derivatives of every order
there. Adding the tensor series preserves this property.
The Euclidean metric itself admits an isometric immersion into
\(\R^3\), whereas the constructed metric does not on any
neighborhood of the origin.
\end{proof}
